%% LyX 2.5.1 created this file.  For more info, see https://www.lyx.org/.
%% Do not edit unless you really know what you are doing.
\documentclass[14pt,twoside,american]{extarticle}
\PassOptionsToPackage{natbib=true}{biblatex}
\usepackage[T1]{fontenc}
\usepackage[utf8]{inputenc}
\usepackage{mathrsfs}
\usepackage{mathtools}
\usepackage{dsfont}
\usepackage{amsmath}
\usepackage{amsthm}
\usepackage{amssymb}
\usepackage{geometry}
\geometry{verbose,tmargin=2.5cm,bmargin=2.5cm,lmargin=2cm,rmargin=3cm}
\usepackage{setspace}
\usepackage{wasysym}
\onehalfspacing

\makeatletter
%%%%%%%%%%%%%%%%%%%%%%%%%%%%%% Textclass specific LaTeX commands.
\numberwithin{equation}{section}
\numberwithin{figure}{section}

%%%%%%%%%%%%%%%%%%%%%%%%%%%%%% User specified LaTeX commands.
\usepackage[utf8]{inputenc}
\usepackage[T1]{fontenc}
% \usepackage[default]{comicneue}

\usepackage{hyperref}
\hypersetup{
	colorlinks=true, %set true if you want colored links
	linktoc=all,     %set to all if you want both sections and subsections linked
	linkcolor=black,  %choose some color if you want links to stand out
}

% \usepackage{url}

%\usepackage[style=alphabetic]{biblatex}
%\addbibresource{bibliography.bib}

\usepackage{graphicx}
\usepackage{enumitem}
\usepackage{booktabs}
\usepackage{csquotes}
% \usepackage{emptypage}
\usepackage{subcaption}
\usepackage{multicol}
\usepackage[dvipsnames]{xcolor}
% \usepackage[usenames, dvipsnames]{xcolor}
\usepackage{amsmath, mathtools, amssymb}
% \usepackage{amsfonts} 
\usepackage{newpxtext} 
\usepackage[euler-digits]{eulervm}
%\usepackage{libertine,libertinust1math}

%\usepackage[sc,osf]{mathpazo}   % With old-style figures and real smallcaps.
%\linespread{1.025}              % Palatino leads a little more leading

% Euler for math and numbers
%\usepackage[euler-digits,small]{eulervm}
\usepackage{extarrows}
\usepackage{amsthm}
\usepackage{mathrsfs}
\usepackage{cancel}
\usepackage{bm}
\usepackage{systeme}
\usepackage{stmaryrd}
\usepackage{dsfont}


\let\implies\Rightarrow
\let\impliedby\Leftarrow
\let\iff\Leftrightarrow
\let\epsilon\varepsilon
\newcommand\contra{\scalebox{1.1}{$\lightning$}}
\let\emptyset\varnothing


\definecolor{correct}{HTML}{009900}
\newcommand\correct[2]{\ensuremath{\:}{\color{red}{#1}}\ensuremath{\to }{\color{correct}{#2}}\ensuremath{\:}}
\newcommand\green[1]{{\color{correct}{#1}}}


\newcommand\hr{
	\noindent\rule[0.5ex]{\linewidth}{0.5pt}
}


% hide parts
\newcommand\hide[1]{}



% si unitx
\usepackage{siunitx}
\sisetup{locale = FR}
% \renewcommand\vec[1]{\mathbf{#1}}
\newcommand\mat[1]{\mathbf{#1}}


% tikz
\usepackage{tikz}
\usepackage{tikz-cd}
\usetikzlibrary{intersections, angles, quotes, calc, positioning}
\usetikzlibrary{arrows.meta}
\usepackage{pgfplots}
\pgfplotsset{compat=1.13}



\tikzset{
	force/.style={thick, {Circle[length=2pt]}-stealth, shorten <=-1pt}
}

% theorems
\makeatother
\usepackage{thmtools}
\usepackage[framemethod=TikZ]{mdframed}
\mdfsetup{skipabove=1em,skipbelow=0em}


% \theoremstyle{plain}

\declaretheoremstyle[
	headfont=\bfseries\sffamily\color{ForestGreen!70!black}, bodyfont=\normalfont,
	mdframed={
			linewidth=2pt,
			rightline=false, topline=false, bottomline=false,
			linecolor=ForestGreen, backgroundcolor=ForestGreen!5,
		}
]{thmgreenbox}

\declaretheoremstyle[
	headfont=\bfseries\sffamily\color{NavyBlue!70!black}, bodyfont=\normalfont,
	mdframed={
			linewidth=2pt,
			rightline=false, topline=false, bottomline=false,
			linecolor=NavyBlue, backgroundcolor=NavyBlue!5,
		}
]{thmbluebox}

\declaretheoremstyle[
	headfont=\bfseries\sffamily\color{RoyalPurple}, bodyfont=\normalfont\itshape,
	mdframed={
			linewidth=2pt,
			rightline=false, topline=false, bottomline=false,
			linecolor=Rhodamine, backgroundcolor=Rhodamine!5,
		}
]{thmpinkbox}

\declaretheoremstyle[
	headfont=\bfseries\sffamily\color{NavyBlue!70!black}, bodyfont=\normalfont,
	mdframed={
			linewidth=2pt,
			rightline=false, topline=false, bottomline=false,
			linecolor=NavyBlue
		}
]{thmblueline}

\declaretheoremstyle[
	headfont=\bfseries\sffamily\color{RawSienna!70!black}, bodyfont=\normalfont\itshape,
	mdframed={
			linewidth=2pt,
			rightline=false, topline=false, bottomline=false,
			linecolor=RawSienna, backgroundcolor=RawSienna!5,
		}
]{thmredbox}

\declaretheoremstyle[
	headfont=\bfseries\sffamily\color{RawSienna!70!black}, bodyfont=\normalfont,
	numbered=no,
	mdframed={
			linewidth=2pt,
			rightline=false, topline=false, bottomline=false,
			linecolor=RawSienna, backgroundcolor=RawSienna!1,
		},
	qed=\qedsymbol
]{thmproofbox}

\declaretheoremstyle[
	headfont=\bfseries\sffamily\color{NavyBlue!70!black}, bodyfont=\normalfont,
	numbered=no,
	mdframed={
			linewidth=2pt,
			rightline=false, topline=false, bottomline=false,
			linecolor=NavyBlue, backgroundcolor=NavyBlue!1,
		},
]{thmexplanationbox}

\declaretheorem[style=thmgreenbox, name=Definition, numberwithin=section]{definition}
\declaretheorem[style=thmbluebox, name=Example, numberwithin=section]{eg}
\declaretheorem[style=thmpinkbox, name=Exercise, numberwithin=section]{ex}
\declaretheorem[style=thmredbox, name=Proposition, numberwithin=section]{prop}
\declaretheorem[style=thmredbox, name=Theorem, numberwithin=section]{theorem}
\declaretheorem[style=thmredbox, name=Lemma, numberwithin=section]{lemma}
\declaretheorem[style=thmredbox, numbered=no, name=Corollary, numberwithin=section]{corollary}

\declaretheorem[style=thmproofbox, name=Proof]{replacementproof}
\renewenvironment{proof}[1][\proofname]{\vspace{-10pt}\begin{replacementproof}}{\end{replacementproof}}


\declaretheorem[style=thmexplanationbox, name=Proof]{tmpexplanation}
\newenvironment{explanation}[1][]{\vspace{-10pt}\begin{tmpexplanation}}{\end{tmpexplanation}}

\declaretheorem[style=thmblueline, numbered=no, name=Remark]{remark}
\declaretheorem[style=thmblueline, numbered=no, name=Note]{note}

\newtheorem*{uovt}{UOVT}
\newtheorem*{notation}{Notation}
\newtheorem*{previouslyseen}{As previously seen}
\newtheorem*{problem}{Problem}
\newtheorem*{observe}{Observe}
\newtheorem*{property}{Property}
\newtheorem*{intuition}{Intuition}


\usepackage{etoolbox}
\AtEndEnvironment{vb}{\null\hfill$\diamond$}%
\AtEndEnvironment{intermezzo}{\null\hfill$\diamond$}%
% \AtEndEnvironment{opmerking}{\null\hfill$\diamond$}%

% http://tex.stackexchange.com/questions/22119/how-can-i-change-the-spacing-before-theorems-with-amsthm
\makeatletter
% \def\thm@space@setup{%
%   \thm@preskip=\parskip \thm@postskip=0pt
% }

\newcommand{\oefening}[1]{%
	\def\@oefening{#1}%
	\subsection*{Oefening #1}
}

\newcommand{\suboefening}[1]{%
	\subsubsection*{Oefening \@oefening.#1}
}

\newcommand{\exercise}[1]{%
	\def\@exercise{#1}%
	\subsection*{Exercise #1}
}

\newcommand{\subexercise}[1]{%
	\subsubsection*{Exercise \@exercise.#1}
}


\usepackage{xifthen}

\def\testdateparts#1{\dateparts#1\relax}
\def\dateparts#1 #2 #3 #4 #5\relax{
	\marginpar{\small\textsf{\mbox{#1 #2 #3 #5}}}
}

\def\@lesson{}%
\newcommand{\lesson}[3]{
	\ifthenelse{\isempty{#3}}{%
		\def\@lesson{Lecture #1}%
	}{%
		\def\@lesson{Lecture #1: #3}%
	}%
	\subsection*{\@lesson}
	\testdateparts{#2}
}

% \renewcommand\date[1]{\marginpar{#1}}


% fancy headers

\usepackage{fancyhdr}
\pagestyle{fancy}
% \fancyhead[LE,RO]{Gilles Castel}
% \fancyhead[RO,LE]{\@lesson}
% \fancyhead[RE,LO]{}
% \fancyfoot[LE,RO]{\thepage}
% \fancyfoot[C]{\leftmark}
% \fancyhf{}
\fancyhf[roh,leh]{\thepage}
\fancyhf[reh]{\nouppercase{\leftmark}}
\fancyhf[loh]{\nouppercase{\rightmark}}
\cfoot{}
\setlength{\headheight}{15pt}
\renewcommand{\headrulewidth}{0pt}
\renewcommand{\footrulewidth}{0pt}


\makeatother




% notes
% \usepackage{todonotes}
\usepackage{tcolorbox}

\tcbuselibrary{breakable}
\newenvironment{verbetering}{\begin{tcolorbox}[
			arc=0mm,
			colback=white,
			colframe=green!60!black,
			title=Opmerking,
			fonttitle=\sffamily,
			breakable
		]}{\end{tcolorbox}}

\newenvironment{noot}[1]{\begin{tcolorbox}[
			arc=0mm,
			colback=white,
			colframe=white!60!black,
			title=#1,
			fonttitle=\sffamily,
			breakable
		]}{\end{tcolorbox}}


\newcommand{\outermeasure}{{{\mu^{\star}}}}

%\def\mbb#1{\mathbb{#1}}
%\def\mfk#1{\mathfrak{#1}}

\newcommand{\inv}{^{\raisebox{.2ex}{$\scriptscriptstyle-1$}}}

\makeatother

\theoremstyle{definition}
\newtheorem*{example*}{\protect\examplename}
\theoremstyle{plain}
\newtheorem{thm}{\protect\theoremname}
\theoremstyle{definition}
\newtheorem{example}[thm]{\protect\examplename}
\newtheorem*{xca*}{\protect\exercisename}
\newtheorem*{sol*}{\protect\solutionname}
\theoremstyle{remark}
\newtheorem*{rem*}{\protect\remarkname}
\usepackage{babel}
\usepackage[style=authoryear,backend=bibtex]{biblatex}
\providecommand{\examplename}{Example}
\providecommand{\exercisename}{Exercise}
\providecommand{\remarkname}{Remark}
\providecommand{\solutionname}{Solution}
\providecommand{\theoremname}{Theorem}

\begin{document}
\title{Aluffi, Chapter-0}
\author{pika}

\maketitle
\global\long\def\bN{\mathbb{N}}%

\global\long\def\bQ{\mathbb{Q}}%

\global\long\def\bZ{\mathbb{Z}}%

\global\long\def\bR{\mathbb{R}}%

\global\long\def\bC{\mathbb{C}}%

\global\long\def\bB{\mathbb{B}}%

\global\long\def\powerset{\mathscr{P}}%

\global\long\def\image{\operatorname\{Im\}}%

\global\long\def\dom{\operatorname{dom}}%

\global\long\def\id{\operatorname\{Id\}}%

\global\long\def\Span{\operatorname\{span\}}%

\global\long\def\one{\mathds{1}}%

\global\long\def\calA{\mathcal{A}}%

\global\long\def\calB{\mathcal{B}}%

\global\long\def\sigmaalgebra{\mathcal{M}}%

\global\long\def\calR{\mathcal{R}}%

\global\long\def\calE{\mathcal{E}}%

\global\long\def\Hom{\mathsf{Hom}}%

\global\long\def\End{\mathsf{End}}%

\global\long\def\Aut{\mathsf{Aut}}%

\global\long\def\catC{\mathsf{C}}%

\global\long\def\Obj{\mathsf{\mathsf{Obj}}}%

\global\long\def\morphism{\longrightarrow}%



\begin{example*}
\citet[Eg-3.4]{aluffi} Define $\catC$ as follows: Let $S$ be any
set, and define a category $\catC\coloneqq\hat{\mathsf{S}}$ by setting
\end{example*}
\begin{itemize}
\item $\Obj(\hat{\mathsf{S}})=\powerset(S)$
\item for any pair of objects $A,B\in\Obj(\hat{\mathsf{S}})$ let $\Hom_{\catC}(A,B)$
be the pair $(A,B)$ if $A\subseteq B,$ and let $\Hom_{\catC}(A,B)=\emptyset$
otherwise. 
\end{itemize}
Indeed, notice that since every $A\in\Obj(\hat{\mathsf{S}})$ satisfies
$A\subseteq A,$ every $\Hom_{\catC}(A,A)$ is nonempty with the identity
morphism $1_{A}$ being the pair $(A,A)$. 

Next, if there is morphism $A\morphism B$ and a morphism $B\to C$
then this means 

that there is a morhpism $A\to C$ which is just the pair $(A,C)$.
Indeed, since $A\subseteq B$ and $B\subseteq C$ implies that $A\subseteq C$,
so $\Hom_{\catC}(A,C)$ is nonempty. 

Next, if there are morphisms $f:A\morphism B$, $g:B\morphism C$
and $h:C\morphism D$ for objects $A,B,C,D\in\Obj(\hat{\mathsf{S}})$
then the morphism $hg=(B,D)$, $f=(A,B),$ $gf=(A,C)$ and so $(hg)f=(A,D)=h(gf)$.
So, associativity also holds. 

That $A\overset{f}{\longrightarrow}B$ satisfies $f1_{A}=f$ and $1_{B}f=f$
is obvious because $1_{A}=(A,A)$ and $f=(A,B)$ so the morphism $f1_{A}$
is just $(A,B)$ while the morphism $1_{B}f$ is also just $(A,B)$. 
\begin{example*}
\citet[Eg-3.3]{aluffi} Fix a set $S$ and a relation $\sim$ that
is reflexive and transitive. Then define a category $\catC$ as such.
\end{example*}
\begin{itemize}
\item the objects of $\catC$ are elements of $S$. 
\item if $a,b\in\Obj(\catC)$ then $\Hom_{\catC}(a,b)$ is a set containing
the element $(a,b)$ if $a\sim b$ and $\Hom_{\catC}(a,b)=\emptyset$
otherwise. 
\end{itemize}
By reflexivity, $\Hom(a,a)$ is always nonempty because it contains
the element $1_{a}=(a,a)$ always. 

If $f:a\to b$ and $g:b\to c$ are any two morphisms then this means
that $a\sim b$ and $b\sim c$. By transitivity, $a\sim c$ and so
$gf:A\to C$ is a well defined morphism. In fact it is the element
$(a,c).$ That is, $\Hom(a,c)=\{gf\}=\{(a,c)\}$. 

If $f:a\to b,g:b\to c$ and $h:c\to d$ are morphisms in $\catC$
then 

$hg=(b,d)$ and $gf=(a,c)$ so $(hg)f=(a,d)=h(gf)$. 

Next, for any morphism $f:a\to b$, $1_{b}f=(a,b)=f$ and $f1_{a}=(a,b)=f$. 

An interesting fact about this category is that given any morphism
$f:a\to b$, take any object $c$ and a pair of morphisms $\alpha',\alpha'':c\to a$.
If $c,\alpha',\alpha''$ are so that
\[
f\alpha'=f\alpha''
\]
 then this implies that $\alpha'=\alpha''$. In fact $\alpha'=\alpha''$
is always true because $\Hom_{\catC}(c,a)$ is always a singleton
or empty. This discussion elucidates the fact that every morphism
in this category is a \emph{monomorphism}. 

In fact, take any morphism $f:a\to b$ and fix any object $c\in\Obj(\catC)$.
Take any pair of morphisms $\beta',\beta'':b\to c$ so that $\beta'f=\beta''f$.
Then, $\beta'=\beta''$. 

Once again, $\beta'=\beta''$ always holds for the same reason as
above. Also, $\beta'f=\beta''f$ also always holds because $\beta'f=(b,c)$
and so is $\beta''f$. Every morphism in this category is an \emph{epimorphism.}
\begin{example*}
\citet[Eg-3.5]{aluffi} Let $\catC$ be a category and fix $A\in\Obj(\catC)$.
We will define another category from $\catC$, denoted $\catC_{A}$
as 
\end{example*}
\begin{itemize}
\item objects of $\catC_{A}$ are morphisms from objects of $\catC$ to
$A$. That is, elements of $\Obj(\catC_{A})$ are of the form $Z\overset{f}{\longrightarrow}A$
for any object $Z\in\Obj(\catC)$. 
\item morphisms between two objects $f,g\in\Obj(\catC_{A})$ are diagrams
comprising of objects of $\catC$ and the appropriate morphisms between
them. 
\end{itemize}
Given two objects $f\colon Z_{1}\to A$ and $g\colon Z_{2}\to A$
in $\catC_{A}$ a morphism between $f,g$ in $\catC_{A}$ is a \emph{commutative
diagram }of the form 

% https://q.uiver.app/#q=WzAsMyxbMSwyLCJDIl0sWzAsMCwiWl8xIl0sWzIsMCwiWl8yIl0sWzIsMCwiZyJdLFsxLDAsImYiLDJdLFsxLDIsIlxcc2lnbWEiLDAseyJzdHlsZSI6eyJib2R5Ijp7Im5hbWUiOiJkYXNoZWQifX19XV0=
\[\begin{tikzcd}
	{Z_1} && {Z_2} \\
	\\
	& A
	\arrow["\sigma", dashed, from=1-1, to=1-3]
	\arrow["f"', from=1-1, to=3-2]
	\arrow["g", from=1-3, to=3-2]
\end{tikzcd}\]

We will show that this construction of $\catC_{A}$ forms a category. 

First, we will show that $\End_{\catC_{A}}(f)$ for any $f\in\Obj(\catC_{A})$
is nonempty and contains the identity morphism by noticing that the
diagram 

% https://q.uiver.app/#q=WzAsMyxbMSwyLCJDIl0sWzAsMCwiWiJdLFsyLDAsIloiXSxbMiwwLCJmIl0sWzEsMCwiZiIsMl0sWzEsMiwiMV9aIiwwLHsic3R5bGUiOnsiYm9keSI6eyJuYW1lIjoiZGFzaGVkIn19fV1d
\[\begin{tikzcd}
	Z && Z \\
	\\
	& A
	\arrow["{1_Z}", dashed, from=1-1, to=1-3]
	\arrow["f"', from=1-1, to=3-2]
	\arrow["f", from=1-3, to=3-2]
\end{tikzcd}\]

where $1_{Z}$ is the identity morphism for the object $Z\in\Obj(\catC)$.
This diagram is the element $1_{f}\in\Hom_{\catC_{A}}(f,f)$. 

Take any two morphisms $f_{1}\to f_{2}$ and $f_{2}\to f_{3}$ in
$\catC_{A}$ next.

We will show that this defines a new morphism which is the composition
of these two morphisms. For this we concatenate the commutative diagrams
$f_{1}\to f_{2}$ and $f_{2}\to f_{3}$ side by side. 

% https://q.uiver.app/#q=WzAsNCxbMiwyLCJBIl0sWzAsMCwiWl8xIl0sWzIsMCwiWl8yIl0sWzQsMCwiWl8zIl0sWzEsMCwiZl8xICIsMl0sWzIsMCwiZl8yIl0sWzEsMiwiXFxzaWdtYSJdLFszLDAsImZfMyJdLFsyLDMsIlxcdGF1Il1d
\[\begin{tikzcd}
	{Z_1} && {Z_2} && {Z_3} \\
	\\
	&& A
	\arrow["\sigma", from=1-1, to=1-3]
	\arrow["{f_1 }"', from=1-1, to=3-3]
	\arrow["\tau", from=1-3, to=1-5]
	\arrow["{f_2}", from=1-3, to=3-3]
	\arrow["{f_3}", from=1-5, to=3-3]
\end{tikzcd}\]

We will show that the following diagram contained in this diagram
also commutes. 

% https://q.uiver.app/#q=WzAsMyxbMiwyLCJBIl0sWzAsMCwiWl8xIl0sWzQsMCwiWl8zIl0sWzIsMCwiZl8zIl0sWzEsMCwiZl8xICIsMl0sWzEsMiwiXFxzaWdtYSBcXHRhdSJdXQ==
\[\begin{tikzcd}
	{Z_1} &&&& {Z_3} \\
	\\
	&& A
	\arrow["{\tau\sigma}", from=1-1, to=1-5]
	\arrow["{f_1 }"', from=1-1, to=3-3]
	\arrow["{f_3}", from=1-5, to=3-3]
\end{tikzcd}\]

Indeed, since $f_{1}=f_{2}\sigma$ and $f_{2}=f_{3}\tau$ from the
previous diagram due to commutativity so $f_{1}=f_{3}\tau\sigma=f_{3}(\tau\sigma)$
which proves that this diagram also commutes, therefore, the composition
of the morphisms $f_{1}\to f_{2}$ and $f_{2}\to f_{3}$ is also a
commutative diagram which describes the morphism $f_{1}\to f_{3}$
between objects $f_{1}$ and $f_{3}$ of the category $\catC_{A}.$

Next, suppose there are three morphisms $f_{1}\to f_{2}\to f_{3}\to f_{4}$
of objects $f_{1},f_{2},f_{3}$ and $f_{4}$ in $\Obj(\catC_{A})$.
That means there are three diagrams as below % https://q.uiver.app/#q=WzAsNSxbMiwyLCJBIl0sWzIsMCwiWl8yIl0sWzAsMCwiWl8xIl0sWzQsMCwiWl8zIl0sWzQsMiwiWl80Il0sWzEsMCwiZl8yICIsMl0sWzIsMCwiZl8xICIsMl0sWzQsMCwiZl80Il0sWzIsMSwiXFxhbHBoYSJdLFszLDAsImZfMyJdLFsxLDMsIlxcYmV0YSJdLFszLDQsIlxcZ2FtbWEiXV0=
\[\begin{tikzcd}
	{Z_1} && {Z_2} && {Z_3} \\
	\\
	&& A && {Z_4}
	\arrow["\alpha", from=1-1, to=1-3]
	\arrow["{f_1 }"', from=1-1, to=3-3]
	\arrow["\beta", from=1-3, to=1-5]
	\arrow["{f_2 }"', from=1-3, to=3-3]
	\arrow["{f_3}", from=1-5, to=3-3]
	\arrow["\gamma", from=1-5, to=3-5]
	\arrow["{f_4}", from=3-5, to=3-3]
\end{tikzcd}\] 

Since $f_{1}\to f_{2}$ is a morphism of objects $f_{1}$ and $f_{2}$
in $\catC_{A}$, the associated diagram above for the two morphisms
is commutative. The commutativity relation is $f_{1}=f_{2}\alpha$.
Similarly, we argue for the other three diagrams and find the following
equalities by commutativity: 

\begin{align*}
f_{2} & =f_{3}\beta\text{ and}\\
f_{3} & =f_{4}\gamma
\end{align*}

Let us set $\sigma\coloneqq f_{1}\to f_{2},\tau\coloneqq f_{2}\to f_{3},\delta\coloneqq f_{3}\to f_{4}$
as our diagrams. Then, 

the diagram $\delta\tau$ is as

% https://q.uiver.app/#q=WzAsMyxbMCwyLCJBIl0sWzAsMCwiWl8yIl0sWzIsMiwiWl80Il0sWzEsMCwiZl8yICIsMl0sWzIsMCwiZl80Il0sWzEsMiwiXFxnYW1tYSBcXGJldGEgIl1d
\[\begin{tikzcd}
	{Z_2} && \\
	\\
	A && {Z_4}
	\arrow["{f_2 }"', from=1-1, to=3-1]
	\arrow["{\gamma \beta }", from=1-1, to=3-3]
	\arrow["{f_4}", from=3-3, to=3-1]
\end{tikzcd}\]

and the diagram $(\delta\tau)\sigma$ is

% https://q.uiver.app/#q=WzAsMyxbMCwwLCJaXzEiXSxbMCwyLCJBIl0sWzIsMiwiWl80Il0sWzAsMiwiZl80IFxcZ2FtbWEgXFxiZXRhIFxcYWxwaGEgIl0sWzIsMSwiZl80ICJdLFswLDEsImZfMSIsMl1d
\[\begin{tikzcd}
	{Z_1} && \\
	\\
	A && {Z_4}
	\arrow["{f_1}"', from=1-1, to=3-1]
	\arrow["{\gamma \beta \alpha }", from=1-1, to=3-3]
	\arrow["{f_4 }", from=3-3, to=3-1]
\end{tikzcd}\]

Next, the diagram $\tau\sigma$ is 

\[\begin{tikzcd}
	{Z_1} && {Z_3}\\
	\\
	A \\
	\arrow["{f_1}"', from=1-1, to=3-1]
	\arrow["{\beta \alpha}", from=1-1, to=1-3]
	\arrow["{f_3 }", from=1-3, to=3-1]
\end{tikzcd}\]

and finally the diagram for $\delta(\tau\sigma)$ is 

\[\begin{tikzcd}
	{Z_1} && \\
	\\
	A && {Z_4}
	\arrow["{f_1}"', from=1-1, to=3-1]
	\arrow["{\gamma \beta \alpha }", from=1-1, to=3-3]
	\arrow["{f_4 }", from=3-3, to=3-1]
\end{tikzcd}\]

So, $\delta(\tau\sigma)=(\delta\tau)\sigma$ proves associativity. 

That given any morphism $\alpha:f_{1}\to f_{2}$ of $\catC_{A}$,
$1_{f_{2}}\alpha=\alpha$ and $\alpha1_{f_{1}}=\alpha$ is obvious
from the same kinds of steps taken above. 
\begin{example*}
\citet[Eg-3.6]{aluffi}
\end{example*}
For the category $\catC$ associated to the set of integers with the
$\leq$ being $\sim$ as in Eg-3.3. 

fix an object $A=3$. Then, any object of $\catC_{A}$ is a pair $(x,3)\in\bZ\times\bZ$
with $x\leq3$. There is a morphism 

$(m,3)\to(n,3)$ if and only if $m\leq n$. 

We learned about the slice category, now let us see an example of
a category that is dual to this - the coslice category.
\begin{example}
\citet[Eg-3.8]{aluffi} For a concrete example, take the category
$\mathsf{Set}$ for which $\Obj(\mathsf{Set)}$ is the class of all
sets. Fix a singleton $\{*\}$ in $\mathsf{Set}$. Next, we define
a new category $\mathsf{Set}^{*}$ whose objects are all the morphisms
$f:\{*\}\to A$ for any object $A$ of $\mathsf{Set}.$ Any object
in $\mathsf{Set}$ may be identified as a pair $(S,s)$ where $s=f(*)$
because the information of any object in $\mathsf{Set}^{*}$ comprises
of the choice of the nonempty set $S$ and of an element $s\in S$
that $f$ maps $*$ to in $S$. 

A morphism between two objects $(S,s)$ and $(T,t)$ in $\mathsf{Set}^{*}$
corresponds to a map $\sigma:S\to T$ \emph{such that }$\sigma(s)=t$. 

Objects in $\mathsf{Set}^{*}$ are called \emph{pointed sets}. 
\end{example}

\begin{xca*}
\citet[Ex-3.1]{aluffi} Let $\catC$ be a category. Consider a structure
$\catC^{\mathsf{op}}$ with 
\end{xca*}
\begin{itemize}
\item $\Obj(\catC^{\mathsf{op}})\coloneqq\Obj(\catC)$;
\item for a pair of objects $A,B\in\Obj(\catC^{\mathsf{op}})$, $\Hom_{\catC^{\mathsf{op}}}(A,B)\coloneqq\Hom_{\catC}(B,A)$ 
\end{itemize}
that is, given any pair of objects $A,B$ of the category $\catC^{\mathsf{op}}$,
arrows from $A$ to $B$ are arrows from $B$ to $A$ in $\catC$. 
\begin{sol*}
Take any pair of objects $A,B\in\Obj(\catC^{\mathsf{op}})$. Define
\[
\ocircle_{\catC^{\mathsf{op}}}\colon\Hom_{\catC^{\mathsf{op}}}(A,B)\times\Hom_{\catC^{\mathsf{op}}}(B,C)\to\Hom_{\catC^{\mathsf{op}}}(A,C)
\]

by the following rule: 

\[
g\ocircle_{\catC^{\mathsf{op}}}f\coloneqq f\ocircle g
\]
 where the $\ocircle$ in the right hand side is the composition of
$f$ and $g$ in $\catC$. 
\end{sol*}
Since $\Hom_{\catC}(A,A)$ is nonempty for every $A\in\Obj(\catC)$
and objects of $\catC^{\mathsf{op}}$ are simply objects in $\catC$,
$\Hom_{\catC^{\mathsf{op}}}(A,A)$ is also nonempty and contains the
same $1_{A}\in\Hom_{\catC}(A,A)$. 

The morphism $f\ocircle_{\catC^{\mathsf{op}}}g$ is an element of
$\Hom_{\catC^{\mathsf{op}}}(A,C)$ because $f\circ g$ is an element
of $\Hom_{\catC}(C,A)$. The well definedness of the map $\ocircle_{\catC^{\mathsf{op}}}$follows
immediately from the well definedness of $\ocircle$. Just take elements
$f,f'\in\Hom_{\catC}(B,A)$ and $g,g'\in\Hom_{\catC}(C,B)$ with $f=f'$,
$g=g'$ and then notice that 
\begin{align*}
g\ocircle_{\catC^{\mathsf{op}}}f= & f\ocircle g\\
= & f'\ocircle g'\\
= & g'\ocircle_{\catC^{\mathsf{op}}}f'
\end{align*}

Next, take the morphisms $f:A\to B,g:B\to C,h:C\to D$ in $\catC^{\mathsf{op}}$. 

Then, 
\begin{align*}
(h\ocircle_{\catC^{\mathsf{op}}}g)\ocircle_{\catC^{\mathsf{op}}}f & =f\ocircle(g\ocircle h)\\
 & =(f\ocircle g)\ocircle h\\
 & =h\ocircle_{\catC^{\mathsf{op}}}\big(g\ocircle_{\catC^{\mathsf{op}}}f\bigr)
\end{align*}

showing that the composition is associative. 

Finally, pick a morphism $f\in\Hom_{\catC^{\mathsf{op}}}(A,B)$. Then,
\begin{align*}
1_{B}\ocircle_{\catC^{\mathsf{op}}}f= & f\ocircle1_{B}\\
= & fp
\end{align*}

and 
\begin{align*}
f\ocircle_{\catC^{\mathsf{op}}}1_{A}= & 1_{A}\ocircle f\\
= & f
\end{align*}

This completes our construction and proof of construction of the opposite
category.
\begin{rem*}
We choose to indicate the morphism $f\in\Hom_{\catC^{\mathsf{op}}}(A,B)$
by just the same symbol when we consider it in $\Hom_{\catC}(B,A)$.
This might not be a bit tricky to deal with initially but it does
make writing the proof quite more convenient. $\hfill\blacksquare$
\end{rem*}
\begin{xca*}
\citet[Ex-3.2]{aluffi} If $A$ is a finite set, how large is $\End_{\mathsf{Set}}(A)$?
\end{xca*}
\begin{sol*}
The morphisms in $\mathsf{Set}$ are simply functions between the
objects (sets) of $\mathsf{Set}$. So we are simply counting the number
of functions from a finite set to itself. 

I claim that for a finite set $A$ of $n$ elements, there are $n^{n}$
functions $A\to A$. Since bijections compose to give bijections,
it suffices to prove the statement for a set of the form $\{1,\dots,n\}$. 

We notice that there is a bijection $\{1,\dots,n\}^{\{1,\dots,n\}}\to\prod^{n}_{k=1}\{1,\dots,n\}$.
Indeed, the map 
\[
f\mapsto(f(i))^{n}_{i=1}
\]

is the obvious candidate. By the product rule of combinatorics, we
see that $\prod^{n}_{k=1}\{1,\dots,n\}$ has cardinality $n^{n}$
and we are done. $\hfill\blacksquare$
\end{sol*}
\begin{xca*}
\citet[Ex-3.3]{aluffi} Formulate precisely what it means to say that
$1_{a}$ is an identity with respect to composition in Eg-3.3. 
\end{xca*}
\begin{sol*}
Take any pair of objects $a,b\in\Obj(\catC)$.If $\Hom_{\catC}(a,b)$
is empty, then the statement is trivially true. If not then pick $(a,b)\in\Hom_{\catC}(a,b)$.
It is clear that we have 
\[
(a,b)1_{a}=(a,b)(a,a)=(a,b)
\]

which means that $1_{a}$ is the identity with respect to composition.
We can do similar for $1_{b}.$ $\hfill\blacksquare$
\end{sol*}
\begin{xca*}
\citet[Ex-3.4]{aluffi} Can we define a category in the style of Eg-3.3
using the relation $<$ on $\bZ.$
\end{xca*}
\begin{sol*}
Every $\Hom_{\catC}(a,a)$ needs to be nonemtpy. If we try to define
a category using < we miss out on $a\sim a$ which does not give us
any obvious candidate of $1_{a}\in\Hom_{\catC}(a,a)$. Shortly put,
$<$ is not reflexive. $\hfill\blacksquare$
\end{sol*}
\begin{xca*}
\citet[Ex-3.5]{aluffi} Define in what sense the category constructed
in Eg-3.4 coincides with that in Eg-3.3.
\end{xca*}
\begin{sol*}
Given a set $S$, define the relation $\sim$ on $\powerset(S)$ as
such: any two subsets $A,B$ of $S$ are related by $\sim$, as in
$A\sim B$ if and only if $A\subseteq B$. The rest is obvious. 

Simply put, the $\subseteq$ relation is a preorder. That is, $A\subseteq A$
for every susbet $A$ of $S$ and $A\subseteq B,B\subseteq C$ means
$A\subseteq C$ for any subsets $A,B,C$ of $S$. $\hfill\blacksquare$
\end{sol*}
\begin{xca*}
\citet[Ex-3.8]{aluffi} A \emph{subcategory} $\catC'$ of a category
$\catC$ consists of a collection of objects of $\catC$, with morphisms
$\Hom_{\catC'}(A,B)\subseteq\Hom_{\catC}(A,B)$ for all objects $A,B\in\Obj(\catC')$,
such that identities and compositions in $\catC$ make $\catC'$ into
a category. A subcategory $\catC'$ is full if $\Hom_{\catC'}(A,B)=\Hom_{\catC}(A,B)$
for all $A,B\in\Obj(\catC')$. Construct a category of \emph{infinite
sets }and explain how it may be viewed as a full subcategory of $\mathsf{Set}.$
\end{xca*}
\begin{sol*}
We define a category $\mathsf{InfSet}$ as such:
\end{sol*}
\begin{itemize}
\item the objects of $\mathsf{InfSet}$ are simply the infinite sets ni
$\Obj(\mathsf{Set)}$, that is $\Obj(\mathsf{InfSet)}\subseteq\Obj(\mathsf{Set})$
\item given any two infinite sets $A,B\in\Obj(\mathsf{InfSet)}$ we define
the set of morphisms $A\to B$ by $\Hom_{\mathsf{InfSet}}(A,B)\coloneqq\Hom_{\mathsf{Set}}(A,B)$. 
\end{itemize}
That these constructions make $\mathsf{InfSet}$ into a subcategory
of $\mathsf{Set}$ is obvious. It has the same identity and compositions
as $\mathsf{Set}.$ 

That it is a full subcategory of $\mathsf{Set}$ is by definition.
$\hfill\blacksquare$
\begin{xca*}
\citet[Ex-3.6]{aluffi} Define a category $\mathsf{V}$ by taking
$\Obj(\mathsf{V})\coloneqq\bN$ and letting $\Hom_{\mathsf{v}}(n,m)\coloneqq\text{the set of }m\times n\text{ matrices with real entries, for all }n,m\in\bN.$
Use product of matrices to define composition. Does this category
\emph{feel }familiar? 
\end{xca*}
\begin{sol*}
Given objects $j,k,l\in\Obj(\mathsf{V)},$ define morphisms $M_{k\times j}\in\Hom_{\mathsf{V}}(j,k),N_{l\times k}\in\Hom_{\mathsf{V}}(k,l)$
where $M_{k\times j}$ and $N_{l\times k}$ are $k\times j$ and $l\times k$
shaped matrices of reals respectively. Define composition

$\Hom(j,k)\times\Hom(k,l)\to\Hom(j,l)$

by the rule 

$M_{k\times j}\circ N_{l\times k}\coloneqq N_{l\times k}M_{k\times j}$ 

where the product on the right hand side is matrix multiplication
which is well defined for $l\times k$ and $k\times j$ matrices.
The result is a matrix $l\times j$ which is an element of $\Hom_{\mathsf{V}}(j,l)$. 

Given any morphism $M_{n\times m}$, consider the morphism $1_{n}\in\Hom(n,n)$.
Notice that $1_{n}M_{n\times m}=M$ and $M_{n\times m}1_{m}=M$ which
shows that $V$ is a category. 

This is very familiar to the category $\mathsf{Vect}$ of all finite
dimensional $k-$vector spaces. The objects are finite dimensional
spaces and the morphisms are linear maps. The spaces can be identified
by their dimension which is a natural number and the linear maps can
be identified by matrices. This suggests there is a sense of \emph{isomorphism
}or \emph{sameness} between these two categories. Such ideas have
not been covered yet in the book so we will not dig further into them
for now. 

The matrix of 0 rows and 0 columns can represent a linear operator
on the trivial vector space $\{0\}$ which has dimension 0. $\hfill\blacksquare$
\end{sol*}
\begin{xca*}
\citet[Ex-3.7]{aluffi} Define carefully objects and morphisms in
Eg-3.7, and draw the diagram corresponding to composition.
\end{xca*}
\begin{sol*}
Suppose $\catC$ is given as a category. Then, fix some object $S\in\Obj(\catC)$.
Define the new category $\catC^{S}$ as follows: 
\end{sol*}
\begin{itemize}
\item objects in $\catC^{S}$ are morphisms $S\to A$ for objects $A\in\Obj(\catC)$, 
\item morphisms between objects $f_{1,}f_{2}$ in $\catC^{s}$ are diagrams
as follows: 
\end{itemize}
\[\begin{tikzcd}
	& S \\ 
	\\
	{Z_1} && {Z_2} \\
	\arrow["\sigma", dashed, from=3-3, to=3-1]
	\arrow["f_{1}"', from=1-2, to=3-1]
	\arrow["f_{2}", from=1-2, to=3-3]
\end{tikzcd}\]

with $f_{1}=\sigma f_{2}$; that is to say that the diagrams are commutative. 

Here, $Z_{1},Z_{2}\in\Obj(\catC)$. 

Given any object $f:S\to Z\in\Hom_{\catC}(S,Z)$ of $\catC^{S}$ where
$Z\in\Obj(\catC)$ there is a diagram 

\[\begin{tikzcd}
	& S \\ 
	\\
	{Z} && {Z} \\
	\arrow["{1_f}", dashed, from=3-1, to=3-3]
	\arrow["{f}"', from=1-2, to=3-1]
	\arrow["{f}", from=1-2, to=3-3]
\end{tikzcd}\]

which obviously commutes and therefore, if we denote this diagram
by $1_{f}$ then $1_{f}\in\Hom_{\catC^{S}}(f,f)$ for every $f\in\Obj(\catC^{S})$. 

Given two morphisms $f_{1}\to f_{2}\to f_{3}$ in $\catC^{S}$, we
will define their composition by the following commutative diagram: 

\[\begin{tikzcd}
	{Z_1} && {Z_2} && {Z_3} \\
	\\
	&& A && \\
	\arrow["{f_1}", from=3-3, to=1-1]
	\arrow["{f_2}", from=3-3, to=1-3]
	\arrow["{f_3}"', from=3-3, to=1-5]
	\arrow["{\alpha}", to=1-1, from=1-3]
	\arrow["{\beta}", to=1-3, from=1-5]
\end{tikzcd}\]

This diagram indeed commutes because $f_{1}=\alpha f_{2},f_{2}=\beta f_{3}$
so $f_{1}=\alpha\beta f_{3}$

which gives the commutative diagram 

\[\begin{tikzcd}
	{Z_1} & &  {Z_3} \\
	\\
	& A & \\
	\arrow["{f_1}", from=3-2, to=1-1]
	\arrow["{f_3}"', from=3-2, to=1-3]
	\arrow["{\alpha\beta}", from=1-3, to=1-1]
\end{tikzcd}\]

which describes a morphism $f_{1}\to f_{3}$ proving that compositions
exist and are well defined. 

That $1_{f}$ is the identity with respect to composition can be proved
from this easily. We shall only prove the associativity of compositions
next. 

For this, take three morphisms $f_{1}\to f_{2}\to f_{3}\to f_{4}$
of objects $f_{1},f_{2},f_{3},f_{4}$ in $\catC^{S}.$ 

Consider the diagram 

\[\begin{tikzcd}
	{Z_1} && {Z_2} && {Z_3} \\
	\\
	&& A && {Z_4}
	\arrow["{f_1}", to=1-1, from=3-3]
	\arrow["{f_2}", to=1-3, from=3-3]
	\arrow["{f_3}", to=1-5, from=3-3]
	\arrow["{f_4}", to=3-5, from=3-3]
	\arrow["{\alpha}"', from=1-3, to=1-1]
	\arrow["{\beta}"', from=1-5, to=1-3]
	\arrow["{\gamma}", from=3-5, to=1-5]
\end{tikzcd}\]

Next, we will label the morphisms $\sigma:f_{3}\to f_{4},\tau:f_{2}\to f_{3},\lambda:f_{1}\to f_{2}$
and look first at $(\sigma\tau)\lambda$. It is clear that the composition
$\sigma\tau$ has the diagram without $Z_{3}$ and $f_{3}$ in the
middle with an arrow $Z_{4}\to Z_{2}$ that is the morphism $\beta\gamma$.
This, composed with $f_{1}$ gives us a diagram with no $Z_{2}$ and
$f_{2}$ in the middle with a map $Z_{4}\to Z_{1}$ given with the
arrow $\alpha\beta\gamma.$ 

Next, we look at $\sigma(\tau\lambda)$. Clearly, $\tau\lambda$ gives
us the diagram without $Z_{2}$ and $f_{2}$ in the middle with a
morphism from $Z_{3}\to Z_{1}$ given as $\alpha\beta$. Composing
this new diagram with the diagram $\sigma$ erases the $f_{3}$ and
$Z_{3}$ in the middle giving an arrow $Z_{4}\to Z_{1}$ which is
the morphism $\alpha\beta\gamma$. This completes the proof. $\hfill\blacksquare$
\begin{xca*}
\citet[Ex-3.9]{aluffi} Define a morphism between multisets, obtaining
a category $\mathsf{MSet}$ containing a copy of $\mathsf{Set}$ as
a full subcategory. Which objects in $\mathsf{MSet}$ determine ordinary
multisets defined in $2.2$ and how? What does a morphsim of multiset
look like here? 
\end{xca*}
\begin{sol*}
TODO 
\end{sol*}
\begin{xca*}
\citet[Ex-3.10]{aluffi} Since the objects of a category $\catC$
are not necessarily sets, it is not clear how to make sense of a notion
of 'subobjects' in general. In some situations, it does make sense
to talk about subobjects, and the subobjects of any given object A
in $\catC$ are in one-to-one correspondence with the morphism $A\to\Omega$
for a fixed, special object $\Omega$ of $\catC$, called a subobject
classifier. Show that $\mathsf{Set}$ has a subobject classifier.
$\hfill\blacksquare$
\end{xca*}
\begin{sol*}
Fix $\Omega\coloneqq\{0,1\}\in\Obj(\mathsf{Set})$. Clearly every
subobject of a set $A$ in $\Obj(\mathsf{Set})$ should be its subset
and every subset $B$ of $A$ can be identified by its indicator $B\to\Omega,b\mapsto1$
if element $b$ of $A$ satisifes $b\in B$ and $b\mapsto0$ otherwise.
The set of subobjects of $A$ should be the $\powerset(A)$ so we
will establish a bijection $f:\powerset(A)\to\{0,1\}^{A}.$

This bijection will be given by the map $B\mapsto\one_{B}$ where
$\one_{B}$ is the indicator of the set $B$ in $A$. This map is
injective because $\one_{C}=\one_{B}$ says an element $x\in A$ satisfies
$x\in C$ if and only if $x\in B$. So, $C=B$. 

For surjectivity, suppose we are given a map $g\in\{0,1\}^{A}$. We
will find a subset $B\subseteq A$ so that $f(B)=g$. Set $B\coloneqq g^{-1}(1)$.
Then, $f(B)=\one_{B}.$ If we show that $\one_{B}=g$ on $A$ then
we are done. So, take any $a\in A$. If $a\in B$ then $\one_{B}(a)=1$
and $g(a)=1$ because $a\in g^{-1}(1)$. If $a\in A\setminus B$ then
$\one_{B}(a)=0$ while $g(a)=0$ because $a\notin g^{-1}(1)$ and
every point of $A$ must map under $g$ to either 0 or 1. Therefore,
$g=\one_{B}$ and so $f(B)=\one_{B}=g$.

This shows that $\{0,1\}$ is a subobject classifier for $\mathsf{Set}.$
\end{sol*}
%
If in a set $S$ we equip it by the equality relation $\sim\coloneqq=$
then the category $\catC$ whose objects are $\Obj(\catC)=S$ and
morphisms are simply identities $(a,a)$ for every object $a\in S$
(all the morphisms are identities because no $a\sim b$ for distinct
$a$ and $b$in $S$). This is a catetgory in which every morphism
is an isomorphism because theyre all identity. $\hfill\blacksquare$
\begin{xca*}
\citet[Ex-4.1]{aluffi} Suppose we are given $n$ different morphisms
$f_{1},\dots,f_{n}$ so that $f_{k}\circ f_{k+1}$ is always well
defined for $k=1,\dots,n-1$. Then, any valid expression involving
parenthesis $($, $)$, elements $f_{1,}\dots,f_{k}$ and the composition
operator $\circ$ is always independent of the placement of the parenthesis. 
\end{xca*}
\begin{sol*}
Let us denote by $E_{k}$ a well formed expression of length $k$.
By this we mean that there are $k$ morphisms $f_{1},\dots,f_{k}$
in $E_{k}$, $k-1$ composition operators $\circ$ and an arbitrary
number of appropriately nested parenthesis. 

It suffices to show that any $E_{k}$ can be equivalently written
in the form $(\dots(f_{1}\circ f_{2})\circ f_{3})\circ f_{4})\circ\dots)\circ f_{k-1})\circ f_{k}$
which allows us to speak of the unique legal expression involving
$k$ morphisms $f_{1},\dots,f_{k}$ and we can simply write that as
$f_{1}\circ f_{2}\circ\dots\circ f_{k}.$ 

We proceed by strong induction on $n$. 

For the base case of $n=3$, $f_{1}\circ(f_{2}\circ f_{3})=(f_{1}\circ f_{2})\circ f_{3}$
is simply the associativity of $\circ$ in $\catC$. 

Our induction hypothesis is that $K_{k}=(\dots(f_{1}\circ f_{2})\circ f_{3})\circ f_{4})\circ\dots)\circ f_{n-1})\circ f_{n}$
for all expressions $K_{k}$ of length $k\in\bZ_{>0}$ with $3\leq k\leq n$. 

Now let $K_{n+1}$ have length $n+1$. Then, there are $l,m\in\bZ_{>0}$
such that $l+m=n+1$ and expressions $K_{l}$ and $K_{m}$ such that
$K_{n+1}=K_{l}\circ K_{m}.$ Now, there are two cases: 

Case 1: $m=1$. Then, $l=n$, $K_{m}=f_{n+1}$. So, by the induction
hypothesis 

\[
K_{l}=(\dots(f_{1}\circ f_{2})\circ\dots)\circ f_{n-1})\circ f_{n}
\]
 Consequently, 
\[
K_{n+1}=(\dots(f_{1}\circ f_{2})\circ\dots)\circ f_{n-1})\circ f_{n})\circ f_{n+1}.
\]

Case 2: m > 1. By the induction hypothesis, $K_{m}$ can be written
in the form $K_{m}=K_{m-1}\circ f_{n+1}$, and so 

\[
K_{n+1}=K_{l}\circ(K_{m-1}\circ f_{n+1})=(K_{l}\circ K_{m-1})\circ f_{n+1}
\]

But, $K_{l}\circ K_{m-1}$ is itself an expression of length $n$,
so, by the induction hypothesis, 
\[
K_{l}\circ K_{m-1}=(\dots(f_{1}\circ f_{2})\circ f_{3})\circ\dots)\circ f_{n}
\]

Finally, 
\[
K_{n+1}=(\dots(f_{1}\circ f_{2})\circ f_{3})\circ\dots)\circ f_{n})\circ f_{n+1}
\]
 which completes the induction step. $\hfill\blacksquare$
\end{sol*}
\begin{xca*}
\citet[Ex-4.2]{aluffi} Given a set $S$ and a relation $\sim$ defined
on $S$, define the category $\catC$ as such: $\Obj(\catC)=S$ and
given objects $a,b\in S,$if $a\sim b$ then there is a morphism $a\overset{f}{\longrightarrow}b\in\Hom_{\catC}(a,b)$
given by $f=(a,b)$ and so $\Hom_{\catC}(a,b)=\{(a,b)\}$ and if $a$
and $b$ are not related by $\sim$ then $\Hom_{\catC}(a,b)=\emptyset$.
For what conditions on $\sim$ is $C$ a groupoid? 
\end{xca*}
\begin{sol*}
We claim that if $\sim$ is an equivalence relation then $\catC$
is a groupoid. Indeed, we saw that $\catC$ is a category if $\sim$
is reflexive and transitive. If $\sim$ is also symmetric, that is,
if for any $a,b\in S,$ $a\sim b$ then $b\sim a$. Indeed in this
case, every morphism $f:a\to b$ has the left inverse $g:b\to a$
given as \textbf{$(b,a)$ }so that \textbf{$gf=(b,a)(a,b)=(a,a)=1_{a}.$
}Similarly, $f:a\to b$ has the right inverse $g=(b,a)$ so that $(a,b)(b,a)=(b,b)=1_{b}.$
$\hfill\blacksquare$
\end{sol*}
\begin{xca*}
\citet[Ex-4.3]{aluffi} Let $A,B$ be objects of a category $\catC$,
and let $f\in\Hom_{\catC}(A,B)$ be a morphism. 
\end{xca*}
\begin{itemize}
\item Prove that if $f$ has a right inverse, then $f$ is an epimorphism. 
\item Find an epimorphism without a right inverse. 
\end{itemize}
\begin{sol*}
Suppose a morphism $f:A\to B$ between objects $A,B\in\Obj(\catC)$
has a right inverse $g:B\to A$ so that $fg=1_{B}.$ Then, Fix an
object $Z$ and a pair of morphisms $\alpha',\alpha'':B\to Z$ in
$\catC$ so that $\alpha'\circ f=\alpha''\circ f$. This implies that
$\alpha'\circ f\circ g=\alpha''\circ f\circ g$ which implies that
$\alpha'\circ1_{B}=\alpha''\circ1_{B}$ and therefore $\alpha'=\alpha''$. 

Just take the morphism $(2,3)$ in the category given in Eg-3.3. Any
right inverse should make it so that $(2,3)\circ g=(3,3)$ so only
$(3,2)$ is a choice but this is not a morphism that makes sense in
this category.$\hfill\blacksquare$
\end{sol*}
\begin{xca*}
\citet[Ex-4.4]{aluffi} Can you define a subcategory $\catC_{\text{nonmono}}$
of $\catC$ by restricting to morphisms that are \emph{not }monomorphisms? 
\end{xca*}
\begin{sol*}
If $\catC$ is the category in Eg-3.3 then every morphism is \emph{monic.}
In that case the set $\Hom_{\catC_{\text{nonmono}}}(a,a)$ for any
$a\in\Obj(\catC)$ is always empty because even $1_{A}$ is monic. 

It might be worth thinking about whether we can do this if we do admit
$1_{A}$ to every $\Hom_{\catC_{\text{nonmono}}}(a,a)$ in $\catC_{\text{nonmono}}$.
Perhaps we might have to look at some examples like $\mathsf{Grp}$
or $\mathsf{CRing}.$ 
\end{sol*}
\begin{xca*}
\citet[Ex-4.5]{aluffi} Give a concrete description of mono and epi
morphisms in the category $\mathsf{MSet}$ you constructed in Ex-3.9.
\end{xca*}
\begin{sol*}
TODO 
\end{sol*}
%
\begin{thebibliography}{Paolo Aluffi(2009)}
\bibitem[Paolo Aluffi(2009)]{aluffi} 

\end{thebibliography}

\end{document}
